\section{Verfahrens-Evaluation}\label{sec:verfahrens-evaluation}


Für die, im \Kapitel{sec:zielsetzung} genannten, Herausforderungen muss eine geeignete Lösung (Framework/Algorithmus) ermittelt werden.
Dazu werden in den nachfolgenden Kapiteln bekannte Verfahren/Varianten von Regel-basierten Systemen und (Bedingungs-basierten) Logiksystemen vorgestellt und nach einer Gegenüberstellung dieser eins ausgewählt.

\subsection{Bekannte Verfahren}\label{subsec:bekannte-verfahren}

\begin{description}
    \item[\Abk{Prolog}] hat die formalisierte Darstellung von Wissen und die Verwendung von logischen Regeln als Grundlage, um dieses Wissen zu verarbeiten und Probleme zu lösen.
    \Abk{Prolog} basiert auf der \Tech{Horn-Klausel}-Semantik, bei der eine Aussage durch eine Menge von Horn-Klauseln dargestellt und von einem \Abk{Prolog}-Interpreter verarbeitet wird.
    Eine Horn-Klausel ist eine logische Formel, die aus einem \Tech{Kopf} und einer Menge von \Tech{Füßen} besteht, wobei der Kopf ein positives Literal ist und die Füße negative Literale sind.
    Eine Aussage ist dann wahr, wenn es eine Übereinstimmung (Unifikation) zwischen dem Kopf und einer der Füße gibt.
    Um die Horn-Klauseln zu verarbeiten, wird ein Übereinstimmung-Algorithmus genutzt.
    Dieser versucht die Variablen in den Horn-Klauseln so zu substituieren, dass der Kopf und einer der Füße übereinstimmen.
    Nur wenn eine solche Substitution möglich ist, gilt die Horn-Klausel als wahr.
    Dabei kommt ein \sog Backtracking-Mechanismus zum Einsatz, bei dem zu einem früheren Punkt im Suchprozess zurückgekehrt wird, um alternative Lösungen zu finden oder wenn keine Übereinstimmung möglich ist.
    Durch das Navigieren durch verschiedene Möglichkeiten kann \Abk{Prolog} komplexe Probleme lösen.
    \Vgl{PROLOG-Apt1997,PROLOG-Bratko2012}

    \item[\Abk{SAT}] ist ein Erfüllbarkeitsproblem der Aussagenlogik, das darin besteht, festzustellen, ob es eine Boolean-Vereinbarung für eine gegebene \Tech{Boolean-Formel} gibt, welche die Formel als wahr evaluiert.
    Eine Boolean-Formel besteht nur aus den Operatoren \textquote{\Tech{und}}, \textquote{\Tech{oder}} und \textquote{\Tech{nicht}}.
    Sie kann als Konjunktion von Disjunktionen von Literalen dargestellt werden, wobei ein Literal entweder eine Variable oder deren Negation ist.
    \Abk{SAT} ist ein NP-vollständiges Problem, was bedeutet, dass es keine effiziente algorithmische Lösung gibt, die für alle Fälle eine korrekte Lösung findet.
    Es gibt jedoch algorithmische Ansätze, die in vielen Fällen effizient arbeiten und erfolgreich eingesetzt werden.
    Nennenswerte und weit verbreitete Ansätze sind der \Abk{DPLL}-Algorithmus und Backtracking-Algorithmen, wie \AbkS{CDCL}{Conflict-Driven Clause Learning}.
    Letzterer vermeidet unnötiges Backtracking und verkürzt damit die Laufzeit des Algorithmus, indem es Konflikt- und Klausel-Lernstrategien nutzt.
    Das \Abk{SAT}-Solver-Verfahren wird häufig mit anderen Techniken, wie \zB Formalisierungs- und Beweisverfahren, kombiniert, sodass Probleme in einer formalen Sprache beschrieben und mögliche Lösungen effizient verifiziert werden können.
    Es ist ein wichtiges und weit verbreitetes (Entscheidungs-)Problem in der theoretischen Informatik und der künstlichen Intelligenz, das viel Anwendung findet und weiterhin intensiv erforscht wird, um effizientere und leistungsfähigere Lösungsmethoden zu entwickeln.
    \Vgl{SAT-Zhang2002,SAT-MarquesSilva2008,SAT-Vizel2015}

    \item[\Abk{ASP}] ist ein deklaratives Problemlösungsparadigma.
    Es basiert auf der Logikprogrammierung (\AbkK{Prolog}), dem Erfüllbarkeitsproblem der Aussagenlogik (\AbkK{SAT}) sowie der Wissensrepräsentation und Argumentation\Englisch{knowledge representation and reasoning}.
    Ein \Abk{ASP}-System besteht aus einer Menge von \Tech{Regeln} und \Tech{Fakten}.
    Die Regeln sind logische Implikationen in der Form \textquote{\textit{wenn A dann B}}, wobei \Tech{A} und \Tech{B} logische \Tech{Aussagen} sind.
    Ein Beispiel wäre: \textquote{\textit{Wenn es regnet, dann nehme ich einen Regenschirm mit.}}.
    Dies kann interpretiert werden als: wenn die Bedingung \textquote{\textit{es regnet}} erfüllt ist, dann ist die Konsequenz \textquote{\textit{ich nehme einen Regenschirm mit}} auch erfüllt.
    Fakten sind einfache logische Aussagen, die bereits als \Tech{wahr} bekannt sind; sie haben keine Bedigung.
    Unbekannte Aussagen werden, gemäß der \AbkS{CWA}{Complete World Assumption}, als \Tech{falsch} gewertet.
    Das Ziel des \Abk{ASP}-Systems ist es, eine Menge von Aussagen zu finden, die, unter den gegebenen Regeln und Fakten, \Tech{wahr} sind.
    Diese Menge von Aussagen wird als \Tech{Answer Set} (auch \Tech{stable models}) bezeichnet.
    Um ein Answer Set zu finden, verwendet das \Abk{ASP}-System (\Tech{Answer Set Solver}) einen Suchalgorithmus, um den Raum möglicher Lösungen zu durchsuchen.
    Es beginnt mit der Menge der Fakten und wendet die Regeln an, um zu sehen, welche anderen Aussagen logisch aus ihnen abgeleitet werden können.
    Wenn ein Widerspruch gefunden wird, geht das System zurück und versucht einen anderen Weg.
    Wenn es keine weiteren Schlussfolgerungen mehr findet, gibt es die Menge der abgeleiteten Aussagen als Answer Set zurück.
    Weitere Eigenschaften von \Abk{ASP} sind \uA:
    \begin{inparaenum}[\bfseries(a)]
        \item Es ist möglich mehrere Lösungen (Answer Sets) zu finden.
        \item Regeln und Fakten können negiert oder als \textquote{unbekannt/unvollständig} gekennzeichnet werden, weshalb \Abk{ASP} als nicht-deterministisch angesehen werden kann.
        \item Regeln können mit Gewichtungen versehen werden, was es ermöglicht unterschiedliche Arten von Unvollständigkeiten zu berücksichtigen.
        \item Die formale Modellierung ist besonders flexibel, weil Regeln dynamisch hinzugefügt und entfernt werden können.
    \end{inparaenum}
    \Vgl{ASP-Eiter2009,ASP-Brewka2011}
\end{description}


\subsection{Gegenüberstellung und Auswahl}\label{subsec:gegenuberstellung-und-auswahl}

In der \Tabelle{tab:verfahrensgegenüberstellung} werden die vorgestellten Verfahren mit den folgenden Kriterien verglichen:
\begin{itemize*}[label={}, itemjoin={,}, itemjoin*={{ sowie}}]
    \item das Prinzip des Verfahrens
    \item die Sprache, welche die grundlegende Syntax und Semantik vorgibt
    \item die Technik der Lösungssuche
    \item mögliche Anwendungsbereiche
    \item Verbreitung (in der Wissenschaft)
    \item Stärken und Schwächen bei der Anwendung.
\end{itemize*}

\hTable{Verfahrensgegenüberstellung}{tab:verfahrensgegenüberstellung}{
    \begin{tabulary}{\linewidth}{@{}CJJJ@{}}
        \toprule
        %% empty
        & \textbf{\Abk{Prolog}}
        & \textbf{\Abk{SAT}}
        & \textbf{\Abk{ASP}}
        \\\midrule
        \vCell{\textbf{Prinzip}}
        & Logikprogrammierung
        & Wahrheitsbelegung für boolesche Variablen finden
        & Lösungen nicht-deterministisch finden, die Bedingungen\Englisch{constraints} erfüllen
        \\\rowcolor{evenRow}
        \vCell{\textbf{Sprache}}
        & Horn-Klauseln
        & Boolesche Formeln
        & Normalform Horn-Klauseln
        \\
        \vCell{\textbf{Technik}}
        & Unifikation, Backtracking
        & Backtracking, \Abk{DPLL}-Algorithmus
        & Stable Model Semantics, Grounder, Solver
        \\\rowcolor{evenRow}
        \vCell{\textbf{Anwendung}}
        & Künstliche Intelligenz, Expertensysteme, natürliche Sprache
        & Logische Schaltungen, Beweisverfahren, Scheduling- und Optimierungsprobleme
        & Bedingungs-basierte Modellierung, kombinatorische Optimierung, Planungsprobleme, Entscheidungsunterstützung
        \\
        \vCell{\textbf{Verbreitung}}
        & 4320 Google Scholar Treffer mit: \Tech{PROLOG ''programming in logic''}
        & 5170 Google Scholar Treffer mit: \Tech{SAT ''boolean satisfiability problem''}
        & 7970 Google Scholar Treffer mit: \Tech{ASP ''answer set programming''}
        \\\rowcolor{evenRow}
        \vCell{\textbf{Stärken}}
        & Natürliche Ausdrucksweise, einfache Handhabung von Unsicherheiten
        & Effiziente Lösungstechniken, geeignet konsistente Lösungen zu finden
        & Präzise und vollständige Repräsentation von Wissen, unterstützt Negation und Unbestimmtheit, komplexe Programmflusskontrolle
        \\
        \vCell{\textbf{Schwächen}\tnote{(1)}\hspace*{1em}}
        & Ineffiziente Lösungstechniken, begrenzte Skalierbarkeit, ungeeignet für komplexe Kontrollstrukturen\tnote{(2)}
        & Schwierigkeiten bei Unsicherheiten und Wissens-Abhängigkeiten, konstante Lösungssuche \ggf ineffizient bei großen Problemen
        & Komplexe Syntax und Semantik (Lernkurve)
        \\\bottomrule
    \end{tabulary}
    \begin{tablenotes}
        \item[(1)] um die Schwächen zu kompensieren, werden diese Verfahren oft mit Anderen kombiniert.
        \item[(2)] wie Schleifen und bedingte Anweisungen.
    \end{tablenotes}
}

Die Verbreitung wird rudimentär über die Anzahl der gefundenen Ergebnisse über Google Scholar bestimmt.
Es wird nach der expliziten (in Hochkommata gestellten) englischen Langform der Abkürzungen \Abk{Prolog}, \Abk{SAT} und \Abk{ASP} gesucht, wobei diese die Kurzform als Präfix hat.
Das Suchformat entspricht damit: \Begriff{<Kurzform> ''<englische Langform>''}.
Besonderes hohe Bedeutung werden den Kriterien Verbreitung, Stärken und Schwächen zugesprochen.
Die Verbreitung stellt einen sehr einfachen Indikator für die Verbreitung des Verfahrens und somit für das Finden möglicher hilfreicher Ressourcen dar.
Dies ist auch in Hinblick der Langlebigkeit der Software relevant.
Die beiden Kriterien Stärken und Schwächen hingegen zeigen \insb die Eignung des Verfahrens zur Formalisierung eines \Abk{PP}-Regelwerks und der Modellierung einer Charakterüberprüfung.

Demnach ist das \textbf{Verfahren \Abk{ASP}} für die Umsetzung \textbf{am besten geeignet} und somit auch zu wählen.
\Abk{ASP} ist ein modernes und verbreitetes Verfahren und erfüllt damit am ehesten den \hyperref[itm:TA.Popularitaet]{vorletzten Punkt der abgeleiteten technischen Anforderungen auf S.\pageref{itm:TA.Popularitaet}}.
Es bringt die notwendige Effizienz bei großen Problemen auf, welches dem umfangreichen \Abk{DSA}-Regelwerk gerecht wird.
Es bringt ebenfalls ein großes Set an Möglichkeiten mit sich, wie \zB Negationen und komplexe Kontroll- und Flussstrukturen.
Des Weiteren erlaubt es bereits vom Prinzip und der Technik her mehrere mögliche Lösungen, was der Realität, also dem Anwendungsbereich, am ehestem entspricht.
Die Schwäche einer steilen Lernkurve, aufgrund der komplexen Syntax und Semantik, kann und wird in der Realisierung durch umfangreiche Dokumentation beachtet und entgegengewirkt.